% ============================================================
% EJEMPLO - LaTeX: documento con formato profesional
% ============================================================
%
% Que es:
%   LaTeX NO es un lenguaje de programacion: es un sistema para
%   escribir documentos. Tu escribes texto + ordenes (comandos),
%   y el produce un PDF con calidad de imprenta.
%
% Como se prueba (NO necesitas instalar nada):
%   - Entra a https://overleaf.com (gratis), crea un proyecto nuevo
%     y pega todo este archivo.
%   - O instala MiKTeX (https://miktex.org) y compila con:
%         pdflatex 09_latex/01_documento.tex
%
% Fijate en: las formulas matematicas se escriben entre $...$
%   y las secciones con \section{...}. LaTeX numera solo.
% ============================================================

\documentclass[11pt, a4paper]{article}

% Paquetes: herramientas extra (idioma, formulas, tablas, etc.)
\usepackage[utf8]{inputenc}
\usepackage[T1]{fontenc}
\usepackage[spanish]{babel}
\usepackage{amsmath}      % formulas matematicas
\usepackage{booktabs}     % tablas bonitas

\title{Informe de ejemplo en LaTeX}
\author{Tu nombre aqui}
\date{\today}

\begin{document}

\maketitle

\begin{abstract}
Este es un documento minimo de ejemplo. Muestra lo basico de LaTeX:
secciones, formulas, tablas y listas. El PDF final sale con formato
profesional, sin apretar botones de formato.
\end{abstract}

\section{Introduccion}

En LaTeX, el texto se escribe normal. Para destacar algo se usa
\textbf{negrita} o \textit{cursiva}. Las secciones se numeran solas:
esto es la seccion 1, y las subsecciones serian 1.1, 1.2, etc.

\section{Formulas}

Las formulas se escriben entre signos de dolar. Por ejemplo, el promedio:
\begin{equation}
    \bar{x} = \frac{1}{n} \sum_{i=1}^{n} x_i
\end{equation}

No hay que hacer nada para numerarla: LaTeX la numera automaticamente.

\section{Tabla}

\begin{table}[h]
    \centering
    \begin{tabular}{lcc}
        \toprule
        Sitio    & Muestras & Peso promedio (g) \\
        \midrule
        Sitio A  & 10       & 11.72 \\
        Sitio B  & 10       & 15.06 \\
        \bottomrule
    \end{tabular}
    \caption{Resumen de las mediciones.}
    \label{tab:resumen}
\end{table}

La Tabla~\ref{tab:resumen} se referencia sola con su numero.

\section{Conclusion}

LaTeX es muy usado en tesis y articulos cientificos porque:
\begin{itemize}
    \item Numera solo las secciones, figuras, tablas y formulas.
    \item Maneja las referencias y la bibliografia automaticamente.
    \item Produce un PDF con acabado profesional.
\end{itemize}

\end{document}
